\PassOptionsToPackage{colorlinks=true,allcolors=black}{hyperref}
% REQ-FILE: The above line, \PassOptionsToPackage{} must be very first line.

\documentclass[11pt]{article}

% Encoding and font setup for cross-platform reproducibility
\usepackage[T1]{fontenc}
\usepackage{lmodern}

% Improve readability for dense theoretical material
\usepackage{setspace}
\onehalfspacing
\usepackage{caption}
\captionsetup[table]{font=small,skip=4pt}

% Mathematical notation and table formatting for formal semantics
\usepackage{amsmath, amssymb}
\usepackage{array}
\usepackage[table]{xcolor}
\usepackage{colortbl}
\usepackage{booktabs}
\usepackage{graphicx}

% Theorem environments (best-practice, widely standard)
\usepackage{amsthm}

% Transformation classifications
\newcommand{\PRS}{\textsf{PRS}}
\newcommand{\BRK}{\textsf{BRK}}
\newcommand{\NEU}{\textsf{NEU}}
\newcommand{\NA}{\textsf{N/A}}

% Identity regimes
\newcommand{\OBL}{\mathrm{OBL}}
\newcommand{\OCC}{\mathrm{OCC}}
\newcommand{\REC}{\mathrm{REC}}

\newcommand{\LOC}{\mathrm{LOC}}
\newcommand{\OBJ}{\mathrm{OBJ}}

\newcommand{\SCOPEE}{\mathrm{SCOPE\text{-}E}}
\newcommand{\SCOPES}{\mathrm{SCOPE\text{-}S}}

\newcommand{\RULEC}{\mathrm{RULE\text{-}C}}
\newcommand{\RULES}{\mathrm{RULE\text{-}S}}

% Formal notation
\newcommand{\Rspace}{\mathcal{R}}
\newcommand{\Fset}{\mathcal{F}}
\newcommand{\Class}{\mathrm{Class}}
\newcommand{\basis}{\mathrm{basis}}
\newcommand{\Substrate}{\mathcal{S}}
\newcommand{\Frameworks}{\mathcal{F}}
\newcommand{\SubstrateRef}{\mathsf{Ref}(\Substrate)}

% ------------------------------------------------------------
% Numbering convention:
% - All theorem-like statements share one counter, numbered by section.
% - Definition/Example/Remark/Note also share the same counter family.
%   (This is the most common convention in math/logic/category theory writing.)
% ------------------------------------------------------------

% Italic body text: results
\theoremstyle{plain}
\newtheorem{theorem}{Theorem}[section]
\newtheorem{lemma}[theorem]{Lemma}
\newtheorem{proposition}[theorem]{Proposition}
\newtheorem{corollary}[theorem]{Corollary}
\newtheorem{conjecture}[theorem]{Conjecture}

% Claim varies; most standard is to share the theorem counter
\newtheorem{claim}[theorem]{Claim}

% Upright body text: definitions and examples
\theoremstyle{definition}
\newtheorem{definition}[theorem]{Definition}
\newtheorem{example}[theorem]{Example}

% Upright body text: remarks/notes
\theoremstyle{remark}
\newtheorem{remark}[theorem]{Remark}
\newtheorem{note}[theorem]{Note}
\newtheorem{convention}[theorem]{Convention}

% Optional: unnumbered versions for front-matter previews/callouts
% (This is also standard when you want a preview without numbering.)
\newtheorem*{theorem*}{Theorem}
\newtheorem*{lemma*}{Lemma}
\newtheorem*{proposition*}{Proposition}
\newtheorem*{corollary*}{Corollary}
\newtheorem*{conjecture*}{Conjecture}
\newtheorem*{claim*}{Claim}
\newtheorem*{definition*}{Definition}
\newtheorem*{example*}{Example}
\newtheorem*{remark*}{Remark}
\newtheorem*{note*}{Note}
\newtheorem*{convention*}{Convention}

% Bibliography management and hyperlink support
\usepackage{url}
\usepackage{natbib}
\bibpunct{(}{)}{;}{a}{}{ }

% load hyperref before bookmark to avoid warnings
\usepackage{hyperref}
\usepackage{bookmark}

% Support multiple author affiliations
\usepackage{authblk}

% Improve typographic quality and line breaking
\usepackage{microtype}

% Keywords macro for structured abstract metadata
\providecommand{\keywords}[1]{\textbf{Keywords: } #1}

% Paragraph spacing prioritizes readability over traditional indentation
\setlength{\parskip}{0.75em}
\setlength{\parindent}{0em}

% Compact list formatting for dense conceptual content
\usepackage{enumitem}
\setlist[itemize]{itemsep=0.2em, topsep=0.2em, parsep=0em, partopsep=0em}

% =============================================================================
% PREFERRED CITATION SPECIFICATION
%
% @techreport{case2026regimes,
%   author      = {Case, Denise M.},
%   title       = {Referential Regimes: Transformation-Invariant Identity for Neutral Substrates},
%   institution = {Northwest Missouri State University},
%   year        = {2026},
%   doi         = {10.48550/arXiv.2601.16152},
%   url         = {https://doi.org/10.48550/arXiv.2601.16152}
% }
% =============================================================================


% ============================================================
\title{Referential Regimes:  \\
Transformation-Invariant Identity \\
for Neutral Substrates}
% ============================================================
\author{Denise M. Case}
\affil{Northwest Missouri State University, Computer Science and Information Systems, Maryville, MO, USA}
\date{} % do not change

\begin{document}

\maketitle
\vspace{-1em}

\begin{abstract}
  Data systems increasingly operate under persistent
  legal, political, and analytic disagreement,
  where no single interpretive authority can be assumed.
  A \emph{neutral substrate} provides stable shared reference
  without requiring agreement about causal or normative interpretation:
  it fixes reference structurally and leaves interpretation to
  extension layers, following the neutrality-by-design constraint
  developed in \emph{Neutral Substrates}~\citep{case2026neutral}.

  This paper derives the identity structure such a substrate requires,
  using a small transformation algebra rather than a list of kinds.
  We treat identity as transformation-invariance:
  a referent is individuated by the admissible operations that preserve it.
  Once the routine operations that record systems perform are fixed,
  including re-expression, annotation, refinement, decomposition,
  reassignment, forking, and provenance extension,
  individuation becomes a structural calculation.

  That calculation exposes a mismatch between reference kinds and
  identity regimes. An accountability substrate must refer to six kinds
  of thing: obligation-bearers, rules, occurrences, scopes, records, and
  plain referents.
  Three of those kinds are assigned more than one admissible basis.
  A scope may be fixed by its extension or by its structure;
  a rule by its content or by its structure;
  a plain referent by its locus or by its object.
  A single routine operation separates each pair:
  decomposition separates the two scope readings,
  refinement the two rule readings, and forking
  the two plain-referent readings.

  Counting reference kinds gives six;
  transformation analysis forces at least nine identity regimes.
  The three additional regimes add no new kind to the inventory.
  They follow from how identity behaves under operations
  the substrate cannot avoid performing, and any attempt to collapse them
  reappears as a hidden regime that persistent disagreement
  cannot hold fixed.

  The result is a conditional lower bound, monotone in the relevant sense:
  a richer transformation basis may leave the bound unchanged
  or force further splits, but it cannot erase a distinction
  already witnessed here. Its content is diagnostic.
  The additional regimes are the distinctions upper-ontology
  implementations may leave to the modeler, and that cannot be recovered
  under persistent disagreement without hidden regimes.
\end{abstract}

\keywords{
  formal ontology;
  identity and persistence;
  neutral substrates;
  extension stability;
  accountability;
  causal commitment;
  normative commitment
}

\input{sections/01_intro}
% !TeX root = ../se200_identity_regimes.tex

\section{Background: Neutral Substrates}
\label{sec:background}

This paper refines the referential-regime component of the neutral-substrate
constraint introduced in SE-100.
In SE-100, a neutral substrate is a shared foundational layer whose neutrality
is guaranteed at design time by restricting substrate-layer commitments to
referential commitments and permitted attribution propositions.
It makes no object-level substrate-layer commitments to causal or normative
propositions.

The present paper analyzes the identity and persistence structure required
for those referential commitments to remain usable as common ground under
persistent interpretive disagreement.
The regimes derived below should therefore be read as refinements of the
referential regime required by SE-100, not as a general-purpose upper ontology.


% !TeX root = ../se200_identity_regimes.tex

% =============================================================================
\section{Formal Setting}
\label{sec:formal_setting}
% =============================================================================

The results in this paper are conditional on the assumptions fixed here.
These assumptions specify the operating conditions under which the
neutral-substrate constraint of Section~\ref{sec:background} is applied.
The goal is to identify the lower-bound identity structure required of a
neutral accountability substrate under those conditions.

% -----------------------------------------------------------------------------
\subsection{Persistent Disagreement}
\label{subsec:persistent_disagreement}
% -----------------------------------------------------------------------------

We assume environments in which disagreement is persistent and extensions are uncoordinated.
Extensions may introduce
incompatible causal accounts, incompatible normative attributions, and
incompatible interpretive classifications.
They cannot be forced to converge.
Substrate revision may not be used to resolve interpretive conflict.
Consequently, substrate-level admissibility must not depend on agreement among extensions.
Neutrality must be maintained without recourse to convergence.

Persistent disagreement is not mere temporary uncertainty.
It is the condition in which multiple admissible frameworks
may remain simultaneously usable while disagreeing about
causation, responsibility, justification, applicability, or normative significance.
A neutral substrate must therefore support shared reference
without embedding any one framework's interpretation as the substrate-level account.

% -----------------------------------------------------------------------------
\subsection{Stable Reference in Strong Form}
\label{subsec:stable_reference_strong}
% -----------------------------------------------------------------------------

Extension-stable reference is inherited from \emph{Neutral Substrates}.
Here we fix the demanding case in which stability must
survive ordinary transformations of records.
Identity conditions must not depend on interpretive convergence.
An entity that must change its ontological status to enter a
new relationship has not been stably referred to.
A stability-critical distinction may not be carried only by a
role, flag, or contextual predicate.
Such devices rely on downstream conventions and can therefore
permit divergent judgments under persistent disagreement.

This is a deliberately strong setting.
Engineering systems often manage identity pressure through
negotiated, local, or context-relative schemes, which may be useful in
practice but are not the case analyzed here.
This paper studies the case in which stable reference must be available
as common ground before interpretive agreement is available.

% -----------------------------------------------------------------------------
\subsection{Minimality}
\label{subsec:minimality}
% -----------------------------------------------------------------------------

We minimize referential structure subject to the above constraints.
Minimality is measured by the number of identity regimes
needed to preserve neutrality and stable reference.
A reduction is admissible only when it preserves the relevant identity distinctions
without relocating them into hidden discriminators.

A reduction counts as genuine only when it does not reintroduce
a collapsed distinction through
a role, flag, contextual predicate, workflow state, schema convention,
or application-level discriminator.
The criterion is analytic.
It identifies the regime structure required under the stated constraints.
It does not recommend use of the core set beyond
the neutral-substrate conditions analyzed here.

% -----------------------------------------------------------------------------
\subsection{Transformation Basis}
\label{subsec:transformation_basis}
% -----------------------------------------------------------------------------

Identity pressure appears when records are transformed.
We therefore fix a finite transformation basis
\[
\Fset =
\{\mathrm{RE},\mathrm{AN},\mathrm{RF},\mathrm{AD},\mathrm{RC},
\mathrm{RA},\mathrm{SU},\mathrm{BF},\mathrm{PV},\mathrm{SE}\}.
\]
The basis is not claimed to be complete for all possible record systems.
It is a representative basis of ordinary operations in
provenance, versioning, institutional recordkeeping, and data-governance practice.
The lower-bound result is conditional on this basis.
Adding further transformations cannot collapse distinctions
already witnessed by this basis.
It may leave the lower bound unchanged or force additional splits.

\begin{definition}[Transformation basis]
\label{def:transformation_basis}
The transformation basis $\Fset$ consists of the following admissible operation families:
\begin{enumerate}[label=(\alph*),nosep]
  \item $\mathrm{RE}$, re-expression: change of representation format while preserving the represented content;
  \item $\mathrm{AN}$, annotation: addition of non-identity-bearing information;
  \item $\mathrm{RF}$, refinement: increase in descriptive, structural, or normative detail;
  \item $\mathrm{AD}$, aggregation or decomposition: restructuring into coarser or finer components;
  \item $\mathrm{RC}$, re-contextualization: change in applicability context;
  \item $\mathrm{RA}$, reassignment: change in association, bearer, or assignment relation;
  \item $\mathrm{SU}$, substitution: replacement of the identity carrier;
  \item $\mathrm{BF}$, branch or fork: divergence into multiple continuations;
  \item $\mathrm{PV}$, provenance extension: addition of historical, derivational, or evidentiary metadata;
  \item $\mathrm{SE}$, state evolution: change in state of an enduring referent over time.
\end{enumerate}
\end{definition}

The transformation families are structural rather than interpretive.
They do not assert causal or normative claims.
They describe operations that may be logged or recognized without
deciding what those operations mean under any particular admissible framework.
Their significance for this paper lies only in their effect on identity.

% -----------------------------------------------------------------------------
\subsection{Classification of Transformations}
\label{subsec:classification_transformations}
% -----------------------------------------------------------------------------

A regime classifies each applicable transformation by how it affects the declared identity basis.
We use three classifications.

\begin{definition}[Transformation classification]
\label{def:transformation_classification}
For a regime $\tau$ and transformation family $f\in\Fset$, the classification
\[
\Class_\tau(f)\in\{\PRS,\BRK,\NEU,\NA\}
\]
has the following meaning:
\begin{enumerate}[label=(\alph*),nosep]
  \item $\PRS$: the transformation is identity-preserving for $\tau$;
  \item $\BRK$: the transformation is identity-breaking for $\tau$;
  \item $\NEU$: the transformation is identity-neutral for $\tau$;
  \item $\NA$: the transformation is inapplicable to $\tau$.
\end{enumerate}
\end{definition}

The distinction between $\PRS$ and $\NEU$ matters.
A $\PRS$ transformation preserves identity in a way that is relevant to the regime's identity basis.
An $\NEU$ transformation does not bear on that identity basis.
Both are non-breaking.
Only $\BRK$ breaks identity.
Thus $\NEU$ does not mean that the operation is ignored by the system.
It means that the operation does not alter identity under the regime being considered.
A fourth value, $\NA$, marks a transformation that cannot arise for the regime's carrier at all,
as distinct from one that arises but leaves identity untouched.

\begin{definition}[Non-breaking transformation]
\label{def:nonbreaking}
A transformation family $f$ is \emph{non-breaking} for regime $\tau$ when
\[
\Class_\tau(f)\in\{\PRS,\NEU\}.
\]
It is \emph{breaking} for $\tau$ when
\[
\Class_\tau(f)=\BRK.
\]
\end{definition}

% -----------------------------------------------------------------------------
\subsection{Identity Regimes}
\label{subsec:identity_regimes}
% -----------------------------------------------------------------------------

The new primitive of this paper is the identity regime.
An identity regime is not an ontological kind.
It is a transformation profile over an identity basis.

\begin{definition}[Identity regime]
\label{def:identity_regime}
An \emph{identity regime} $\tau$ consists of an identity basis
together with a classification function
\[
\Class_\tau:\Fset\to\{\PRS,\BRK,\NEU,\NA\}.
\]
The identity basis states what counts as sameness for the referent.
The classification function states how each transformation family affects sameness under that basis.
\end{definition}

\begin{definition}[Induced identity relation]
\label{def:induced_identity}
Let $\tau$ be an identity regime.
The induced identity relation $\sim_\tau$ is the equivalence relation
generated by transformations that are non-breaking for $\tau$.
Thus $r\sim_\tau s$ means that $r$ and $s$ are representations of the same referent under regime $\tau$.
\end{definition}

This definition makes identity transformation-relative.
The same pair of representations may coincide under one regime
and diverge under another.
The substrate must declare which identity basis is being preserved.

% -----------------------------------------------------------------------------
\subsection{Hidden Regimes}
\label{subsec:hidden_regimes}
% -----------------------------------------------------------------------------

A hidden regime is a collapsed identity distinction recovered outside
the declared substrate structure.
Hidden regimes obstruct genuine minimality.
They allow a system to appear to use fewer identity regimes
while relying on downstream discriminators to perform the missing identity work.

\begin{definition}[Hidden regime]
\label{def:hidden_regime}
A \emph{hidden regime} is a stability-critical distinction that
is not represented as an identity regime in the substrate,
but is instead recovered through a role, flag, contextual predicate,
workflow state, schema convention, or application-level discriminator.
\end{definition}

Hidden regimes violate the neutral-substrate discipline because their
application is not fixed by the substrate.
They require agreement about how the discriminator should be interpreted.
Under persistent disagreement, that agreement cannot be assumed.
A collapse that can be repaired only by a hidden regime is therefore
not a genuine reduction in declared referential structure.

% -----------------------------------------------------------------------------
\subsection{Scope of the Lower Bound}
\label{subsec:lower_bound_scope}
% -----------------------------------------------------------------------------

The lower bound proved below is conditional.
It holds for neutral accountability substrates satisfying the assumptions
of this section and using the transformation basis of Definition~\ref{def:transformation_basis}.
It asserts that, under the stated constraints,
fewer than nine regimes cannot preserve stable reference
without changing transformation behavior or reintroducing a hidden regime.
The claim is a lower bound only:
it does not assert that nine regimes
are sufficient for every substrate,
nor that the nine derived below exhaust
the coherent identity regimes.

The next section states the reference inventory.
The inventory supplies the referents that must remain addressable.
The transformation basis supplied here determines
which of those referents can share an identity regime.

% !TeX root = ../se200_referential_regimes.tex

% =============================================================================
\section{Reference Inventory}
\label{sec:inventory}
% =============================================================================

Before analyzing identity, we fix the carrier inventory.
An accountability substrate must let one state, inspect, and compare claims
about who is responsible, which authority governs, what happened,
where it applies, what was recorded, and what was acted upon.
It must do this without embedding causal or normative content at the substrate level.

The inventory is modest.
It is not a theorem, an upper ontology, or a complete domain ontology.
It states the carrier kinds required by the accountability setting analyzed here.
Those carriers are the referents that must remain addressable under the assumptions
of Section~\ref{sec:formal_setting}.

The inventory supplies identity carriers but does not determine identity regimes.
A reference kind identifies the referents that must remain addressable.
An identity regime specifies sameness for a carrier under transformation.
The lower-bound proof depends on that distinction.

% -----------------------------------------------------------------------------
\subsection{Reference Kinds and Identity Carriers}
\label{subsec:reference_kinds_identity_carriers}
% -----------------------------------------------------------------------------

\begin{definition}[Reference kind]
  \label{def:reference_kind}
  A \emph{reference kind} is a class of referents that a
  neutral accountability substrate must be able to identify, track, and
  distinguish from other referents without making object-level causal or normative commitments.
\end{definition}

\begin{definition}[Identity carrier]
  \label{def:identity_carrier}
  An \emph{identity carrier} is a referential position in the substrate to
  which an identity regime may be assigned.
  It is the object whose sameness, difference, and persistence are evaluated under transformation.
\end{definition}

Reference kinds supply identity carriers.
They do not themselves determine identity regimes.
A single reference kind may support more than one identity regime when
different identity bases respond differently to the same transformation.
Conversely, the fact that two carriers are both represented in the
substrate does not imply that they share an identity regime.
Regime assignment depends on transformation behavior, not on the name of the kind.

The inventory used in this paper is:

\begin{center}
  \begin{tabular}{@{}l>{\raggedright\arraybackslash}p{0.68\textwidth}@{}}
    \toprule
    Reference kind    & Required referent                                                             \\
    \midrule
    Obligation-bearer &
    a party that can bear responsibility across the occurrences
    discharging, violating, or contesting that responsibility                                         \\
    Rule              &
    a governing normative structure, distinct from the acts that enact, amend, apply, or challenge it \\
    Occurrence        &
    a happening individuated by temporal realization and provenance                                   \\
    Scope             &
    the applicability domain for a rule, including places, parties, periods,
    procedures, classifications, or conditions                                                        \\
    Record            &
    a descriptive indicator or record that carries reference without
    asserting causal or normative truth                                                               \\
    Plain referent    &
    a thing acted upon, regulated, located, transferred, or observed
    that neither bears obligations nor grounds authority as such                                      \\
    \bottomrule
  \end{tabular}
\end{center}

These six reference kinds are familiar from foundational ontology,
legal ontology, provenance, and institutional recordkeeping.
The paper does not claim novelty for the inventory itself.
The inventory matters because it supplies the carriers on which
the later identity analysis operates.

% -----------------------------------------------------------------------------
\subsection{Obligation-Bearers}
\label{subsec:inventory_obligation_bearers}
% -----------------------------------------------------------------------------

An accountability substrate must support reference to obligation-bearing entities.
An obligation-bearer is a party that can bear
responsibility, duty, liability, authority, or accountability
across occurrences in which those statuses are fulfilled, violated, transferred, contested, or discharged.
The substrate need not decide whether a party is responsible in a particular case.
It must preserve reference to the party that a framework may treat as responsible.

In multi-jurisdiction regulatory reporting, a corporate entity may be treated differently across jurisdictions.
One jurisdiction may individuate responsibility at the parent-entity level.
Another may individuate it at the subsidiary, facility, role, or delegated-agent level.
A neutral substrate need not settle which legal theory governs.
It must still allow the relevant bearer to remain referable across those competing frameworks.

Obligation-bearers cannot be reduced to the occurrences in which obligations are discharged.
The same bearer may participate in many occurrences.
The same occurrence may involve several bearers.
Collapsing bearer into occurrence destroys the ability to compare responsibility across time and across frameworks.

% -----------------------------------------------------------------------------
\subsection{Rules}
\label{subsec:inventory_rules}
% -----------------------------------------------------------------------------

An accountability substrate must support reference to governing rules or normative structures.
A rule is a structure that may be
enacted, amended, delegated, interpreted, applied, violated, or repealed.
The substrate need not assert that the rule is valid, justified, binding, or correctly interpreted.
It must preserve reference to the rule as the object of such attributed claims.

In legislative history tracking, a statute must remain referable
independently of the legislative act that enacted it, the amendment
that altered it, the agency action that applied it, and the judicial decision that interpreted it.
A statute can be contested while still remaining the same referent
for purposes of citation, comparison, and provenance.
If the rule is collapsed into its enactment or application occurrence,
then later disagreement about the occurrence changes reference to the rule itself.

Rules also preview one of the split cases below.
A rule may be individuated by imposed requirements
or by textual or internal structure.
That difference is not settled by naming the kind ``rule.''
It is settled by transformation behavior.

% -----------------------------------------------------------------------------
\subsection{Occurrences}
\label{subsec:inventory_occurrences}
% -----------------------------------------------------------------------------

An accountability substrate must support reference to \emph{occurrences}.
An occurrence is a happening individuated by temporal realization and provenance.
Illustrative examples include filings, inspections, votes, transfers, transactions,
observations, decisions, enactments, amendments, and enforcement actions.
The substrate must preserve reference to the occurrence as something that happened,
independently of how competing frameworks evaluate it.

In multi-jurisdiction regulatory reporting, two jurisdictions may disagree about
whether a filing constitutes compliance.
They must nevertheless be able to refer to the same filing.
In legislative history tracking, parties may disagree about the legal significance of an enactment.
They must nevertheless be able to refer to the same enactment occurrence.
In provenance analysis, institutions may disagree about the meaning or reliability of a measurement.
They must nevertheless be able to refer to the same measurement occurrence.

Occurrences cannot be reduced to obligation-bearers, rules, records, or scopes.
They are time-indexed referents.
They may instantiate, affect, or be described by other carriers, but they are not identical with them.

% -----------------------------------------------------------------------------
\subsection{Scopes}
\label{subsec:inventory_scopes}
% -----------------------------------------------------------------------------

An accountability substrate must support reference to scopes.
A scope is the domain where, to whom, or under what conditions a rule applies.
Scopes may be territorial, institutional, temporal, procedural, classificatory, jurisdictional, or conditional.
They may be nested, overlapping, partially coincident, or contested.

In multi-jurisdiction regulatory reporting, the same rule may apply to
different entities, places, facilities, transactions, or reporting periods under different frameworks.
If scope is represented as an attribute of the rule, then a change in
applicability can appear to change the identity of the rule.
If scope is represented as a property of obligation-bearers, then the
domain of applicability cannot be compared independently of the parties within it.
A neutral substrate needs scope as a separate referent.

Scopes also preview one of the split cases below.
A scope may be individuated extensionally, by covered cases,
or structurally, by internal jurisdictional or applicability organization.
Decomposition can preserve one basis while changing the other.
The reference kind ``scope'' therefore does not determine an identity regime.

% -----------------------------------------------------------------------------
\subsection{Records}
\label{subsec:inventory_records}
% -----------------------------------------------------------------------------

An accountability substrate must support reference to descriptive records.
A record is a descriptive indicator, document, dataset entry, observation record,
provenance item, or other referent that can be cited, transferred, compared, annotated, or analyzed without the substrate asserting a causal or normative claim.
The substrate may record that a source asserts a claim.
It does not thereby assert the claim.

In cross-institutional provenance, a record may be produced by one institution,
received by another, annotated by a third, and reanalyzed by a fourth.
The institutions may disagree about what the record shows or whether it supports a conclusion.
They must nevertheless be able to refer to the same record through custody, annotation, and reuse.

Records cannot be reduced to the entities, occurrences, rules, or claims they describe.
Stable reference to records is required independently of those referents,
including when a record supports an attributed causal or normative claim.

% -----------------------------------------------------------------------------
\subsection{Plain Referents}
\label{subsec:inventory_plain_referents}
% -----------------------------------------------------------------------------

An accountability substrate must also support reference to plain referents.
A plain referent is a thing acted upon, regulated, located, transferred, observed,
measured, or described that neither bears obligations nor grounds authority.
Examples include facilities, parcels, instruments, devices, environmental sites,
physical objects, infrastructure components, samples, and assets.
The category is intentionally modest.
It captures referents that are important to accountability without already being
obligation-bearers, rules, occurrences, scopes, or records.

Plain referents are needed because many accountability questions concern things
that are acted upon or regulated.
A regulated facility may be inspected.
A parcel may fall within a jurisdictional scope.
An instrument may produce a measurement record.
A sample may be transferred across institutions.
The substrate must preserve reference to the thing without deciding which
causal or normative interpretation applies to it.

Plain referents also preview one of the split cases below.
A referent may be individuated by locus,
such as the site or place that remains fixed.
It may also be individuated by the object occupying that locus.
Forking or replacement can separate those bases.
The reference kind ``plain referent'' therefore
does not by itself determine an identity regime.

% -----------------------------------------------------------------------------
\subsection{Summary}
\label{subsec:inventory_summary}
% -----------------------------------------------------------------------------

The inventory supplies six reference kinds:
obligation-bearer, rule, occurrence, scope, record, and plain referent.
These kinds are the carriers for the lower-bound construction
and do not determine identity regimes.

The inventory identifies the required referents.
The transformation analysis determines their behavior under ordinary record operations.
The core is derived after applying those operations
to the carrier-basis pairs.

The next section applies the transformation basis of
Section~\ref{sec:formal_setting} to these six carriers.

% !TeX root = ../se200_referential_regimes.tex

% =============================================================================
\section{Transformation-Invariant Identity}
\label{sec:algebra}
% =============================================================================

The previous section fixed the referents the substrate must preserve.
This section analyzes the conditions under which two representations
of such a referent count as the same entity.
The analysis treats identity as \emph{transformation-invariance}:
an identity regime is determined by the behavior of a declared identity basis
under the transformation basis fixed in Section~\ref{sec:formal_setting}.

The governing idea is familiar from formal ontology:
categories are individuated by identity conditions, not by classificatory convenience,
and properties that a thing can gain or lose cannot carry identity~\citep{guarino1998roles,guarinowelty2002}.
This paper applies that idea at the substrate level and in operational form.
A referent remains the same under a regime when the transformations applied
to its representations do not break the identity basis for that regime.
Two representations may be the same under one regime and different under another.
A substrate that leaves the regime implicit has not fixed reference.
It has deferred the identity question to a later interpretive choice.
Under persistent disagreement, that choice cannot be assumed shared.

% -----------------------------------------------------------------------------
\subsection{Identity Bases}
\label{subsec:identity_bases}
% -----------------------------------------------------------------------------

\begin{definition}[Identity basis]
\label{se200.def.IdentityBasis}
An \emph{identity basis} is the feature, structure, role in reference, or persistence condition
whose preservation constitutes remaining the same referent under a given identity regime.
\end{definition}

An identity basis is the condition a regime tracks through transformation.
The lower-bound construction uses responsibility-bearing identity
for obligation-bearers,
temporal-realization and provenance identity for occurrences,
and descriptive-record identity for records.
It also uses three witnessed split cases.
A rule may be content-fixed or structure-fixed.
A scope may be extension-fixed or structure-fixed.
A plain referent may be locus-fixed or object-fixed.

The inventory supplies carriers.
The identity basis determines sameness for a carrier.
The transformation classification determines whether that basis is preserved or broken.
The lower-bound argument begins when two bases for the same carrier kind
classify an ordinary transformation differently.

% -----------------------------------------------------------------------------
\subsection{Profiles and Collapse}
\label{subsec:profiles_and_collapse}
% -----------------------------------------------------------------------------

The formal setting already defined identity regimes as classification profiles over the transformation basis.
We now use the induced identity relation $\sim_\tau$ to compare regimes.

\begin{definition}[Collapse]
\label{se200.def.Collapse}
Two identity regimes $\tau$ and $\tau'$ \emph{collapse} when they induce the same identity relation:
\[
\sim_\tau \ =\ \sim_{\tau'} .
\]
They are \emph{distinct} when some pair of representations is identified under one regime and separated under the other.
\end{definition}

\begin{definition}[Underdetermined kind]
\label{se200.def.underdetermined_kind}
A reference kind is \emph{underdetermined} when its kind label
does not determine a single identity basis and admissible identity bases
for that kind classify some transformation family differently.
Equivalently, one admissible basis for the kind
classifies a transformation as non-breaking,
while another admissible basis for the same kind
classifies that transformation as breaking.
\end{definition}

Underdetermination is the engine of the result.
When basis plurality is in play for an underdetermined kind,
no single regime profile serves all required bases.
The kind must be represented through distinct identity regimes for those bases.
The split is forced by transformation behavior, not by the name of the kind.

% -----------------------------------------------------------------------------
\subsection{Carrier Kinds Not Split in the Core Construction}
\label{subsec:core_unsplit_kinds}
% -----------------------------------------------------------------------------

Three carrier kinds contribute one regime each to the core.
They are not split by the witness operations used in this construction:
obligation-bearers, occurrences, and records.

An \emph{obligation-bearer} is individuated by persistence as the party that can bear
responsibility, duty, liability, authority, or accountability.
Re-expression, annotation, re-contextualization, reassignment, branching,
provenance extension, and state evolution do not alter that identity basis.
Refinement and decomposition preserve identity when they add detail or decompose descriptions
of the bearer without replacing the bearer.
Substitution breaks identity because it replaces the bearer.

An \emph{occurrence} is individuated by temporal realization and provenance.
Re-expression, annotation, refinement, re-contextualization, and provenance extension do not replace the occurrence.
State evolution preserves identity when it tracks phases or stages of the same occurrence.
Decomposition breaks identity when it produces sub-occurrences with their own temporal individuation.
Substitution breaks identity because it replaces the occurrence.
Branching and reassignment are inapplicable to a fixed realization event.

An \emph{record} is individuated by descriptive persistence without substrate-level causal or normative commitment.
Annotation, refinement, decomposition, branching, and provenance extension
preserve record identity because the record survives enrichment, derivation, copying, and descriptive elaboration.
Re-expression, re-contextualization, reassignment, and state evolution are
identity-neutral for the record as such.
Substitution breaks identity because it replaces the record.

These three carrier-basis pairs are included once in the core.
We name their regimes:
\[
\OBL,\qquad \OCC,\qquad \REC .
\]

% -----------------------------------------------------------------------------
\subsection{Witnessed Split Kinds}
\label{subsec:witnessed_split_kinds}
% -----------------------------------------------------------------------------

Three reference kinds have witnessed split pairs in the lower-bound construction.
Each admits at least the two admissible readings used here,
and an ordinary transformation separates those readings.
The pattern is the same in all three cases.
A kind admits two bases $B_1$ and $B_2$.
A transformation $f$ is non-breaking for $B_1$ and breaking for $B_2$.
When admissible frameworks require both bases,
the kind cannot be governed by a single identity regime.

\begin{lemma}[Split pattern]
\label{se200.lem.SplitPattern}
Let $K$ be a reference kind with two admissible identity bases $B_1$ and $B_2$.
Suppose there is a transformation family $f\in\Fset$ and representations $r,s$
such that $f$ preserves identity under $B_1$ but breaks identity under $B_2$.
If a neutral substrate must support both bases for $K$,
then $K$ requires at least two distinct identity regimes.
\end{lemma}

\begin{proof}
If one regime governed both bases, it would assign a single classification to $f$
for the kind.
But the $B_1$ reading requires $f$ to be non-breaking,
while the $B_2$ reading requires $f$ to be breaking.
No single classification can satisfy both readings.
Thus the kind is underdetermined at $f$.

Applying $f$ to the witnessing representations identifies them under $B_1$
and separates them under $B_2$.
The two bases therefore induce distinct identity relations.
When both bases must be supported, the substrate requires at least two
distinct identity regimes for $K$.
\end{proof}

% -----------------------------------------------------------------------------
\subsection{Scope Under Decomposition}
\label{subsec:scope_split}
% -----------------------------------------------------------------------------

A scope may be individuated by its extension or by its structure.
An extension-fixed scope is individuated by the
cases, persons, places, entities, or conditions it covers.
A structure-fixed scope is individuated
by the internal organization through which that coverage is represented.

\begin{proposition}[Scope Underdetermination]
\label{se200.prop.ScopeUnderdetermination}
The scope reference kind is underdetermined at aggregation/decomposition $\mathrm{AD}$.
\end{proposition}

\begin{proof}
Let $r$ be a scope fixed by extension.
Let $s$ be a decomposition of $r$ into subscopes whose union covers the same cases.
The extension is preserved, so $\mathrm{AD}$ is non-breaking for the extension-fixed reading.
Let $r'$ be a scope fixed by structure.
Let $s'$ be a decomposition of $r'$ into a finer internal partition.
The structure is changed, so $\mathrm{AD}$ is breaking for the structure-fixed reading.
Both readings are admissible scopes.
No single classification of $\mathrm{AD}$ can serve both.
By Lemma~\ref{se200.lem.SplitPattern}, when both bases must be supported,
scope requires at least two identity regimes.
\end{proof}

This yields:
\[
\SCOPEE,\qquad \SCOPES .
\]
$\SCOPEE$ is the extension-fixed scope regime.
$\SCOPES$ is the structure-fixed scope regime.
After decomposition, two representations
can apply to the same cases while no longer being the same scope structure.
State evolution is inapplicable to both scope regimes:
a scope is an applicability domain,
not an enduring referent that passes through internal states.

% -----------------------------------------------------------------------------
\subsection{Rule Under Refinement}
\label{subsec:rule_split}
% -----------------------------------------------------------------------------

A rule may be individuated by its content or by its structure.
A content-fixed rule is individuated by its requirements,
permissions, prohibitions, authorizations, or definitions.
A structure-fixed rule is individuated by the organization of the
text, sections, clauses, delegations, or internal rule structure
through which that content is expressed.

\begin{proposition}[Rule Underdetermination]
\label{se200.prop.RuleUnderdetermination}
The rule reference kind is underdetermined at refinement $\mathrm{RF}$.
\end{proposition}

\begin{proof}
Let $r$ be a rule fixed by normative content.
Let $s$ be a refinement of $r$ that preserves the rule's requirements
while adding detail or reorganizing expression.
The content is preserved, so $\mathrm{RF}$ is non-breaking for the content-fixed reading.
Let $r'$ be a rule fixed by textual or internal structure.
Let $s'$ be a refinement of $r'$ that changes that structure.
The structure is changed, so $\mathrm{RF}$ is breaking for the structure-fixed reading.
Both readings are admissible rules.
No single classification of $\mathrm{RF}$ can serve both.
By Lemma~\ref{se200.lem.SplitPattern}, when both bases must be supported,
rule requires at least two identity regimes.
\end{proof}

This yields:
\[
\RULEC,\qquad \RULES .
\]
$\RULEC$ is the content-fixed rule regime.
$\RULES$ is the structure-fixed rule regime.
After refinement, two representations
can require the same thing while no longer being the same text or rule structure.
State evolution is inapplicable to both rule regimes:
a rule is a normative structure, not an enduring referent that passes through internal states.
Changes to rule content are amendments or refinements,
not state transitions of a persisting object.

% -----------------------------------------------------------------------------
\subsection{Plain Referent Under Forking}
\label{subsec:plain_referent_split}
% -----------------------------------------------------------------------------

A plain referent may be individuated by locus or by object.
A locus-fixed referent is individuated by the
place, site, slot, location, or institutional position that remains fixed.
An object-fixed referent is individuated by the
object, asset, instrument, sample, device, parcel, or thing occupying that locus.

\begin{proposition}[Plain-referent Underdetermination]
\label{se200.prop.PlainReferentUnderdetermination}
The plain-referent kind is underdetermined at branch/fork $\mathrm{BF}$.
\end{proposition}

\begin{proof}
Let $r$ be a plain referent fixed by locus.
Let $s$ be a forked continuation over the same locus.
The locus is preserved, so $\mathrm{BF}$ is non-breaking for the locus-fixed reading.
Let $r'$ be a plain referent fixed by object.
Let $s'$ be a fork into non-identical object continuations.
The object identity is not preserved across both continuations,
so $\mathrm{BF}$ is breaking for the object-fixed reading.
Both readings are admissible plain referents.
No single classification of $\mathrm{BF}$ can serve both.
By Lemma~\ref{se200.lem.SplitPattern}, when both bases must be supported,
plain referents require at least two identity regimes.
\end{proof}

This yields:
\[
\LOC,\qquad \OBJ .
\]
$\LOC$ is the locus-fixed plain-referent regime.
$\OBJ$ is the object-fixed plain-referent regime.
After a fork, two representations
can be the same place while no longer being the same object.
State evolution applies to both plain-referent regimes and preserves identity:
a locus or an object is an enduring referent that persists through changes in its state.

% -----------------------------------------------------------------------------
\subsection{Result of the Algebra}
\label{subsec:algebra_result}
% -----------------------------------------------------------------------------

The reference inventory supplied six carrier kinds.
Three carrier-basis pairs contribute one regime each:
\[
\OBL,\qquad \OCC,\qquad \REC .
\]
Three carrier kinds have witnessed split pairs:
\[
\LOC,\ \OBJ,\qquad
\SCOPEE,\ \SCOPES,\qquad
\RULEC,\ \RULES .
\]
Thus the transformation analysis yields the nine-regime lower-bound core:
\[
\begin{aligned}
&\OBL,\quad \OCC,\quad \REC,\quad \LOC,\quad \OBJ,\\
&\SCOPEE,\quad \SCOPES,\quad \RULEC,\quad \RULES .
\end{aligned}
\]

The additional regimes arise within carrier kinds already named.
They are identity bases separated by ordinary transformations.
This is why counting carrier kinds undercounts the structure needed
for stable reference.
The next section collects the nine core regimes in a single classification table.

% !TeX root = ../se200_referential_regimes.tex

% =============================================================================
\section{The Nine-Regime Lower-Bound Core}
\label{sec:nine}
% =============================================================================

The transformation analysis of Section~\ref{sec:algebra}
yields the following nine core regimes.
Three carrier-basis pairs contribute one regime each.
Three carrier kinds have witnessed split pairs.

The table uses the following short regime labels:

\begin{description}[style=nextline,leftmargin=2.5cm]
  \item[$\OBL$] obligation-bearer identity.
  \item[$\OCC$] occurrence identity.
  \item[$\REC$] record identity.
  \item[$\LOC$] locus-fixed plain-referent identity.
  \item[$\OBJ$] object-fixed plain-referent identity.
  \item[$\SCOPEE$] extension-fixed scope identity.
  \item[$\SCOPES$] structure-fixed scope identity.
  \item[$\RULEC$] content-fixed rule identity.
  \item[$\RULES$] structure-fixed rule identity.
\end{description}

Their classifications over $\Fset$ are collected in Table~\ref{tab:master}.
The table records the core profiles used in the lower-bound proof.
Three same-carrier cells carry the split arguments:
BF for $\LOC/\OBJ$,
AD for $\SCOPEE/\SCOPES$,
and RF for $\RULEC/\RULES$.
Together with carrier disjointness, those cells carry the lower-bound result.
The remaining entries summarize diagnostic behavior of the core profiles.

\begin{table}[ht]
  \centering
  \small
  \setlength{\tabcolsep}{4pt}
  \begin{tabular}{@{}l ccc cc cc cc@{}}
    \toprule
       & OBL  & OCC  & REC  & LOC  & OBJ  & SCOPE-E & SCOPE-S & RULE-C & RULE-S \\
    \midrule
    RE & \NEU & \NEU & \NEU & \NEU & \NEU & \NEU    & \NEU    & \NEU   & \NEU   \\
    AN & \NEU & \NEU & \PRS & \NEU & \NEU & \NEU    & \NEU    & \NEU   & \NEU   \\
    RF & \PRS & \PRS & \PRS & \PRS & \PRS & \PRS    & \BRK    & \PRS   & \BRK   \\
    AD & \PRS & \BRK & \PRS & \PRS & \PRS & \PRS    & \BRK    & \PRS   & \BRK   \\
    RC & \NEU & \NEU & \NEU & \NEU & \NEU & \PRS    & \BRK    & \NEU   & \NEU   \\
    RA & \NEU & \NA  & \NEU & \NA  & \NA  & \NEU    & \BRK    & \NEU   & \NEU   \\
    SU & \BRK & \BRK & \BRK & \BRK & \BRK & \BRK    & \BRK    & \BRK   & \BRK   \\
    BF & \NEU & \NA  & \PRS & \PRS & \BRK & \PRS    & \BRK    & \BRK   & \BRK   \\
    PV & \NEU & \NEU & \PRS & \NEU & \NEU & \NEU    & \NEU    & \NEU   & \NEU   \\
    SE & \NEU & \PRS & \NEU & \PRS & \PRS & \NA     & \NA     & \NA    & \NA    \\
    \bottomrule
  \end{tabular}
  \caption{Classification of each regime profile over the transformation basis $\Fset$.
    \PRS{}: preserves identity;
    \BRK{}: breaks identity;
    \NEU{}: identity-neutral/non-breaking;
    \NA{}: not applicable.
    Each witnessed split pair has a canonical separating operation:
    plain referent at BF, scope at AD, and rule at RF.
    The remaining classifications are diagnostic rather than load-bearing
    for the lower-bound proof.}
  \label{tab:master}
\end{table}

The table should be read as a compact statement of identity behavior.
A \PRS{} entry means that the transformation \emph{preserves} the regime's identity basis.
A \BRK{} entry means that the transformation \emph{breaks} that basis and produces a distinct referent under the regime.
An \NEU{} entry means that the transformation is \emph{identity-neutral} for that regime.
An \NA{} entry means that the transformation is \emph{not applicable} to that regime.

Two columns bracket the profiles.
Re-expression is identity-neutral for every core regime:
a change of representation format preserves the declared identity basis.
Substitution is identity-breaking for every core regime:
replacement of the identity carrier breaks identity under each basis.
These uniform columns serve as classification sanity checks.
They set uniform boundary cases for the profile table.
The lower-bound witnesses are the separating cells for the three split pairs.

The witnessed split pairs carry the lower-bound argument.
$\LOC$ and $\OBJ$ are separated by branching or forking.
$\SCOPEE$ and $\SCOPES$ are separated by aggregation or decomposition.
$\RULEC$ and $\RULES$ are separated by refinement.
The canonical separating operation explains why each split is forced.
The two profiles in a split pair may also differ at other transformations.
For scope,
re-contextualization and reassignment are two such non-canonical differences.
They preserve covered cases and are non-breaking for $\SCOPEE$.
They change applicability context or assignment structure and are breaking for $\SCOPES$.

These regimes are defined by identity basis and transformation behavior.
They refine identity bases within the existing carrier kinds.

\begin{definition}[Separation witness]
  \label{se200.def.SeparationWitness}
  A \emph{separation witness} for identity regimes $\tau$ and $\tau'$ is a pair of representations $(r,s)$ such that
  \[
    r \sim_\tau s
    \qquad\text{and}\qquad
    r \not\sim_{\tau'} s,
  \]
  or conversely.
\end{definition}

\begin{proposition}[Distinct core profiles]
  \label{se200.prop.DistinctCoreProfiles}
  The nine profiles of Table~\ref{tab:master} are pairwise distinct.
\end{proposition}

\begin{proof}
  Same-carrier pairs are separated by the split propositions of Section~\ref{sec:algebra}.
  By Proposition~\ref{se200.prop.PlainReferentUnderdetermination}, $\LOC$ and $\OBJ$ are separated at BF.
  By Proposition~\ref{se200.prop.ScopeUnderdetermination}, $\SCOPEE$ and $\SCOPES$ are separated at AD.
  By Proposition~\ref{se200.prop.RuleUnderdetermination}, $\RULEC$ and $\RULES$ are separated at RF.
  In each case there is a separation witness: a pair identified under one regime and separated under the other.

  Profiles with different carriers govern disjoint sets of representations.
  A representation of an obligation-bearer is never a representation of an occurrence,
  an occurrence never a record,
  a rule never a scope,
  a scope never a plain referent,
  and a record never the thing it records.
  Two regimes whose carriers have disjoint representation domains induce
  identity relations that differ on those domains,
  and are therefore distinct by Definition~\ref{se200.def.Collapse}.
  Identifying such profiles would collapse the reference inventory of
  Section~\ref{sec:inventory} and force one carrier to stand in for another.

  Therefore every pair of core regimes in Table~\ref{tab:master}
  is separated either by a same-carrier transformation witness
  or by carrier difference.
  Appendix~\ref{app:sep} records the separation basis for all $\binom{9}{2}=36$ pairs.
\end{proof}

% !TeX root = ../se200_referential_regimes.tex

% =============================================================================
\section{The Lower Bound}
\label{sec:bound}
% =============================================================================

The previous section collected the nine core identity regimes.
This section proves the lower-bound claim.
The proof requires two facts:
the inventory requires stable reference to six carrier kinds,
and three of those carrier kinds exhibit witnessed basis plurality under the transformation basis.

% -----------------------------------------------------------------------------
\subsection{The Inventory Floor}
\label{subsec:inventory_floor}
% -----------------------------------------------------------------------------

The reference inventory of Section~\ref{sec:inventory} requires stable reference to six kinds of carrier:
obligation-bearers, rules, occurrences, scopes, records, and plain referents.
Stable reference to a carrier kind requires an identity regime for that kind,
because the substrate must determine when representations of that carrier count as the same referent.

\begin{lemma}[Inventory floor]
\label{se200.lem.InventoryFloor}
Under the assumptions of Section~\ref{sec:formal_setting},
a neutral accountability substrate realizing the inventory of Section~\ref{sec:inventory}
must realize at least one identity regime for each required carrier kind.
\end{lemma}

\begin{proof}
By Definition~\ref{se200.def.ReferenceKind}, each reference kind in the inventory
must be identifiable, trackable, and distinguishable from other referents
without making object-level causal or normative commitments.
By Definition~\ref{se200.def.IdentityRegime}, an identity regime specifies
sameness for a referent under transformation.
Thus each required carrier kind needs identity conditions sufficient to
preserve stable reference to carriers of that kind.

The required carrier kinds have disjoint representation domains in the inventory.
A representation of an obligation-bearer is not a representation of an occurrence;
an occurrence is not a record;
a rule is not a scope;
a scope is not a plain referent;
and a record is not the thing it records.
Stable reference to each required carrier kind therefore requires identity
conditions for that carrier kind.
Collapsing carrier kinds would collapse the reference inventory.
Therefore each required carrier kind contributes at least one identity regime.
\end{proof}

The inventory therefore gives a floor of six regimes.
The witnessed basis pluralities for plain referents, scopes, and rules
raise that floor to nine.

% -----------------------------------------------------------------------------
\subsection{Witnessed Splits}
\label{subsec:witnessed_splits}
% -----------------------------------------------------------------------------

Three carrier kinds exhibit witnessed basis plurality under the transformation basis.

First, plain referents split into locus-fixed and object-fixed identity.
By Proposition~\ref{se200.prop.PlainReferentUnderdetermination},
branch/fork separates these bases.
A fork may preserve the locus while breaking object identity.
When both bases are in play,
plain referents require both $\LOC$ and $\OBJ$.

Second, scopes split into extension-fixed and structure-fixed identity.
By Proposition~\ref{se200.prop.ScopeUnderdetermination},
aggregation/decomposition separates these bases.
A decomposition may preserve the covered cases while changing the internal scope structure.
When both bases are in play,
scopes require both $\SCOPEE$ and $\SCOPES$.

Third, rules split into content-fixed and structure-fixed identity.
By Proposition~\ref{se200.prop.RuleUnderdetermination},
refinement separates these bases.
A refinement may preserve rule content while changing textual or internal structure.
When both bases are in play,
rules require both $\RULEC$ and $\RULES$.

Each witnessed split adds one regime beyond the one-regime-per-carrier floor.
The additional regimes are not new reference kinds.
They are distinct identity bases for required carrier kinds.

% -----------------------------------------------------------------------------
\subsection{The Bound}
\label{subsec:the_bound}
% -----------------------------------------------------------------------------

\begin{theorem}[Nine-regime lower bound]
\label{se200.thm.NineRegimeLowerBound}
Under the assumptions of Section~\ref{sec:formal_setting},
a neutral accountability substrate that supports the inventory of
Section~\ref{sec:inventory}
and supports the witnessed basis pluralities for plain referents,
scopes, and rules
must realize at least nine identity regimes.
\end{theorem}

\begin{proof}
By Lemma~\ref{se200.lem.InventoryFloor},
the six required carrier kinds in Section~\ref{sec:inventory}
require at least six identity regimes.

Three of those carrier kinds require an additional regime when the witnessed
basis pluralities are in play.
Plain referents require both $\LOC$ and $\OBJ$ by
Proposition~\ref{se200.prop.PlainReferentUnderdetermination}.
Scopes require both $\SCOPEE$ and $\SCOPES$ by
Proposition~\ref{se200.prop.ScopeUnderdetermination}.
Rules require both $\RULEC$ and $\RULES$ by
Proposition~\ref{se200.prop.RuleUnderdetermination}.

Therefore the required core contains:
\[
\begin{gathered}
\OBL,\quad \OCC,\quad \REC,\quad \LOC,\quad \OBJ,\\
\SCOPEE,\quad \SCOPES,\quad \RULEC,\quad \RULES .
\end{gathered}
\]
By Proposition~\ref{se200.prop.DistinctCoreProfiles},
these nine core profiles are pairwise distinct.
Thus the required regimes number at least,
\[
1 + 1 + 1 + 2 + 2 + 2 = 9.
\]
Thus any neutral accountability substrate satisfying the stated assumptions
and supporting the witnessed basis pluralities
must realize at least nine identity regimes.
This proves the lower-bound claim.
\end{proof}

% -----------------------------------------------------------------------------
\subsection{Monotonicity}
\label{subsec:monotonicity}
% -----------------------------------------------------------------------------

\begin{corollary}[Monotonicity]
\label{se200.cor.Monotonicity}
Enlarging the transformation basis cannot lower the bound
established by Theorem~\ref{se200.thm.NineRegimeLowerBound}.
\end{corollary}

\begin{proof}
The core is witnessed by distinctions
already forced by the transformation basis $\Fset$.
Adding transformations adds classification coordinates.
It may leave the existing distinctions unchanged.
It may separate an already-distinct pair in further ways.
It may also force further splits.
But it cannot erase a split already witnessed by $\Fset$.
Therefore any transformation basis containing $\Fset$ still requires
at least the nine core regimes derived here.
\end{proof}

The monotonicity claim is intentionally limited.
It establishes that incompleteness of the transformation basis
cannot make the present lower bound disappear.
A richer basis may leave the bound unchanged or raise it,
but cannot lower it.

% -----------------------------------------------------------------------------
\subsection{Boundary of the Proof}
\label{subsec:bound_not_claim}
% -----------------------------------------------------------------------------

The claim proved here is a lower bound.
Under the stated assumptions and witnessed basis pluralities,
fewer than nine regimes cannot preserve stable reference
without changing transformation behavior or reintroducing a hidden regime.

The proof does not establish an upper bound.
It does not claim that the nine regimes exhaust all coherent identity regimes.
It does not claim that every deployed system must expose these regimes.
Those broader limits are stated in Section~\ref{sec:limits}.

% !TeX root = ../se200_referential_regimes.tex

% =============================================================================
\section{Determinacy}
\label{sec:determinacy}
% =============================================================================

The lower-bound result requires a fixed assignment
from core carrier-basis pairs to identity regimes.
This section records that assignment and
only addresses the carrier-basis pairs used in the lower-bound construction.

Once a core carrier-basis pair is fixed,
the regime assignment is determined.
Which identity basis a real-world referent ought to use requires modeling judgment.

% -----------------------------------------------------------------------------
\subsection{Core Carrier-Basis Pairs}
\label{subsec:core_carrier_basis_pairs}
% -----------------------------------------------------------------------------

\begin{definition}[Core carrier-basis pair]
\label{def:core_carrier_basis_pair}
A \emph{core carrier-basis pair} is one of the carrier kind and identity basis
combinations used in the lower-bound construction.
The core pairs are:
\[
\begin{array}{rcl}
\text{obligation-bearer with responsibility-bearing identity} &\mapsto& \OBL,\\
\text{occurrence with realization-and-provenance identity} &\mapsto& \OCC,\\
\text{record with descriptive-record identity} &\mapsto& \REC,\\
\text{plain referent with locus-fixed identity} &\mapsto& \LOC,\\
\text{plain referent with object-fixed identity} &\mapsto& \OBJ,\\
\text{scope with extension-fixed identity} &\mapsto& \SCOPEE,\\
\text{scope with structure-fixed identity} &\mapsto& \SCOPES,\\
\text{rule with content-fixed identity} &\mapsto& \RULEC,\\
\text{rule with structure-fixed identity} &\mapsto& \RULES.
\end{array}
\]
\end{definition}

The first three pairs contribute one regime each.
The remaining six pairs arise from the witnessed split cases proved in Section~\ref{sec:algebra}.

\begin{definition}[Core regime assignment]
\label{def:core_regime_assignment}
The \emph{core regime assignment} maps each core carrier-basis pair
to the identity regime displayed in Definition~\ref{def:core_carrier_basis_pair}.
\end{definition}

% -----------------------------------------------------------------------------
\subsection{Existence}
\label{subsec:determinacy_existence}
% -----------------------------------------------------------------------------

\begin{lemma}[Existence for the core]
\label{se200.lem.CoreExistence}
Every core carrier-basis pair is assigned to at least one of the nine core regimes.
\end{lemma}

\begin{proof}
\sloppy
The three one-basis core pairs are assigned by construction:
obligation-bearer with responsibility-bearing identity maps to $\OBL$,
occurrence with realization-and-provenance identity maps to $\OCC$,
and record with descriptive-record identity maps to $\REC$.

The remaining core pairs are assigned by the split results of Section~\ref{sec:algebra}.
A plain referent with locus-fixed identity maps to $\LOC$,
and a plain referent with object-fixed identity maps to $\OBJ$,
by Proposition~\ref{se200.prop.PlainReferentUnderdetermination}.
A scope with extension-fixed identity maps to $\SCOPEE$,
and a scope with structure-fixed identity maps to $\SCOPES$,
by Proposition~\ref{se200.prop.ScopeUnderdetermination}.
A rule with content-fixed identity maps to $\RULEC$,
and a rule with structure-fixed identity maps to $\RULES$,
by Proposition~\ref{se200.prop.RuleUnderdetermination}.
Therefore every core carrier-basis pair is assigned to at least one of the nine core regimes.
\end{proof}

% -----------------------------------------------------------------------------
\subsection{Uniqueness}
\label{subsec:determinacy_uniqueness}
% -----------------------------------------------------------------------------

\begin{lemma}[Uniqueness for the core]
\label{se200.lem.CoreUniqueness}
No core carrier-basis pair is assigned to more than one of the nine core regimes.
\end{lemma}

\begin{proof}
Each core carrier-basis pair contains a fixed carrier kind and a fixed identity basis.
The core regime assignment maps that pair
to the regime that preserves the declared basis
under the transformation behavior used in the lower-bound construction.

The three one-basis core pairs map to $\OBL$, $\OCC$, and $\REC$ respectively.
For the split cases, the declared basis selects the side of the split:
locus-fixed plain referent maps to $\LOC$ and not $\OBJ$;
object-fixed plain referent maps to $\OBJ$ and not $\LOC$;
extension-fixed scope maps to $\SCOPEE$ and not $\SCOPES$;
structure-fixed scope maps to $\SCOPES$ and not $\SCOPEE$;
content-fixed rule maps to $\RULEC$ and not $\RULES$;
and structure-fixed rule maps to $\RULES$ and not $\RULEC$.
Thus the assignment is functional on the core pairs.
Therefore no core carrier-basis pair is assigned to more than one core regime.
\end{proof}

% -----------------------------------------------------------------------------
\subsection{Determinacy Theorem}
\label{subsec:determinacy_theorem}
% -----------------------------------------------------------------------------

\begin{theorem}[Core regime assignment]
\label{se200.thm.CoreRegimeAssignment}
Every core carrier-basis pair is assigned to exactly one of the nine core identity regimes.
\end{theorem}

\begin{proof}
\sloppy
Existence follows from Lemma~\ref{se200.lem.CoreExistence}.
Uniqueness follows from Lemma~\ref{se200.lem.CoreUniqueness}.
Therefore every core carrier-basis pair is assigned to exactly one of the nine core regimes.
\end{proof}

This determines the core assignment.
It does not decide whether a modeler has chosen the correct identity basis
for a particular real-world referent.
It says only that the carrier-basis pairs used in the lower-bound construction
have fixed regime assignments.

% -----------------------------------------------------------------------------
\subsection{Failure of Basis Declaration}
\label{subsec:basis_declaration_failure}
% -----------------------------------------------------------------------------

For the witnessed split cases, omission of the basis is not harmless.
A scope without an extension-fixed or structure-fixed basis does not specify
whether decomposition preserves covered cases or internal structure.
A rule without a content-fixed or structure-fixed basis does not specify
whether refinement preserves rule content or textual structure.
A plain referent without a locus-fixed or object-fixed basis does not specify
whether forking preserves location or object identity.
A declared identity basis is required for deterministic assignment.

This observation anticipates the conformance analysis
of the companion \emph{Accountable Records} paper.
In a deployed record system, basis declaration becomes a schema obligation.
In this paper, it marks the boundary of the determinacy result:
regime assignment is fixed for the carrier-basis pairs included
in the core.

% !TeX root = ../se200_identity_regimes.tex

% =============================================================================
\section{The Nine Regimes at Work}
\label{sec:application}
% =============================================================================

The preceding sections derive a nine-regime lower-bound core.
This section shows how they appear in the running diagnostic settings.
The examples are illustrative.
They show that the surplus regimes in the lower-bound core are load-bearing:
they mark distinctions that ordinary transformations force
and that accountability systems must track.

% -----------------------------------------------------------------------------
\subsection{Multi-Jurisdiction Regulatory Reporting}
\label{subsec:app_regulatory}
% -----------------------------------------------------------------------------

A corporation is an obligation-bearer.
It persists across jurisdictional differences in corporate personhood, subsidiary responsibility,
delegated liability, and compliance theory.
It is therefore represented under $\OBL$.

The reporting statute is a rule.
If it is identified by what it requires, then it is represented under $\RULEC$.
A later re-codification that preserves the requirements is the same rule under $\RULEC$
because refinement preserves content-fixed identity.
If the statute is identified by its section structure, then it is represented under $\RULES$.
The same re-codification is then a different rule,
because refinement changes the structure-fixed identity basis.

Each filing, inspection, transaction, or enforcement action is an occurrence.
It is individuated by temporal realization and provenance.
It is therefore represented under $\OCC$.
Two jurisdictions may disagree about whether the occurrence constitutes compliance.
They must still be able to refer to the same occurrence.

The jurisdiction is a scope.
If it is identified by the set of regulated cases it covers, then it is represented under $\SCOPEE$.
Splitting the jurisdiction into districts that together cover the same cases preserves identity under $\SCOPEE$.
If it is identified by its named administrative partition, then it is represented under $\SCOPES$.
The same split changes the structure and therefore breaks identity under $\SCOPES$.

Compliance indicators, filings as documents, measurement reports, and audit records are descriptive records.
They are represented under $\REC$.
They can be compared, annotated, and transferred without the substrate asserting whether compliance actually occurred.

A regulated substance, facility, instrument, parcel, or asset is a plain referent.
If the relevant identity basis is the regulated locus, then it is represented under $\LOC$.
If the relevant identity basis is the object occupying or moving through that locus,
then it is represented under $\OBJ$.

The case exercises all nine regimes in the lower-bound core.
Two jurisdictions can refer to the same corporation, statute, filing, jurisdiction, record,
and facility while evaluating compliance incompatibly.
They can also disagree about whether decomposing a jurisdiction changed the jurisdiction.
That disagreement is the $\SCOPEE/\SCOPES$ distinction.
A kind-count that stops at ``scope'' would miss it.

% -----------------------------------------------------------------------------
\subsection{Legislative History Tracking}
\label{subsec:app_legislative}
% -----------------------------------------------------------------------------

A statute is a rule.
If tracked by normative content, it is represented under $\RULEC$.
An amendment or refinement that preserves what the statute requires preserves identity under $\RULEC$.
If tracked by text, section structure, or internal organization, the statute is represented under $\RULES$.
The same amendment or refinement may then produce a new structure-fixed rule.

The enactment, amendment, vote, administrative application, judicial decision,
or repeal event is an occurrence.
It is represented under $\OCC$.
The occurrence may enact, modify, interpret, or challenge the statute,
but it is not identical with the statute.

The legislature, agency, court, officer, or regulated party is an obligation-bearer
when represented as a party that can bear responsibility, authority, or accountability.
Such entities are represented under $\OBL$.
Their identity does not depend on any single enactment or application event.

Geographic, subject-matter, procedural, temporal, or institutional boundaries are scopes.
A statutory scope may be tracked extensionally by the cases it covers, yielding $\SCOPEE$.
It may also be tracked structurally by the named jurisdictional or procedural partition, yielding $\SCOPES$.
Redistricting, reclassification, or subdivision can preserve coverage while changing structure.

Legislative records, committee reports, interpretive memoranda, docket entries, and compliance histories are records.
They are represented under $\REC$.
They may record that someone asserted a causal or normative claim without
making that claim at the substrate level.

Objects acted upon by the legal regime, such as land parcels, facilities, permits,
instruments, samples, or regulated assets, are plain referents.
They may be locus-fixed under $\LOC$ or object-fixed under $\OBJ$ depending on
what identity basis the substrate declares.

The case shows why the rule split matters.
Legal practice often needs both readings.
A doctrine may remain unchanged while the text is superseded.
A text may remain identifiable while its content is interpreted differently.
The $\RULEC/\RULES$ split prevents these two identity bases from being collapsed.

% -----------------------------------------------------------------------------
\subsection{Cross-Institutional Provenance}
\label{subsec:app_provenance}
% -----------------------------------------------------------------------------

Cross-institutional provenance separates records from the things those records describe.
A dataset entry, measurement record, transfer log, annotation record, or derived metric
is represented under $\REC$.
It may be copied, annotated, refined, and extended with provenance while remaining the same record.

The sampled object, monitored site, instrument, facility, specimen, or collection point
is a plain referent.
If the identity basis is the site, location, or collection point,
then it is represented under $\LOC$.
A fork over records concerning the same site preserves the locus-fixed referent.
If the identity basis is the specimen, instrument, or physical object,
then it is represented under $\OBJ$.
A fork that produces distinct object continuations breaks object-fixed identity.

Producing institutions, custodians, laboratories, agencies, and receiving organizations
are obligation-bearers when represented as parties responsible for custody, transfer, use, retention, or analysis.
They are represented under $\OBL$.

Data-use policies, retention rules, consent terms, sharing agreements, and access restrictions are rules.
When tracked by what they require, they are represented under $\RULEC$.
When tracked by their text, version, or internal structure, they are represented under $\RULES$.

Transfers, annotations, re-analyses, approvals, denials, measurements, and custody changes are occurrences.
They are represented under $\OCC$.
They remain referable as the same events even when institutions disagree about their significance.

Institutional coverage conditions, consent scopes, policy domains, project boundaries,
and data-sharing jurisdictions are scopes.
They may be extension-fixed under $\SCOPEE$ or structure-fixed under $\SCOPES$.
A policy domain may cover the same records after decomposition while no longer having
the same internal structure.

The provenance chain can therefore remain stable while institutions disagree about
classification, causal interpretation, data quality, responsibility, or permissible use.
The $\LOC/\OBJ$ distinction keeps ``the same site'' and ``the same specimen''
from being conflated across a fork.
The record-versus-referent distinction keeps a record from being confused with the thing it records.

% -----------------------------------------------------------------------------
\subsection{Summary}
\label{subsec:app_summary}
% -----------------------------------------------------------------------------

The proof of the lower bound was given in Section~\ref{sec:bound}.
The running cases illustrate why the lower-bound core is not artificial.

Each case uses the three core regimes that are not split in this construction:
$\OBL$, $\OCC$, and $\REC$.
Each case also exercises the three witnessed split pairs:
$\LOC/\OBJ$, $\SCOPEE/\SCOPES$, and $\RULEC/\RULES$.
The surplus regimes correspond to identity distinctions
that ordinary record operations already force.

A jurisdiction that stays the same under decomposition and one that does not are both scopes.
A statute individuated by what it requires and a statute individuated by its text are both rules.
A site and the object occupying it may both be plain referents.
The kind names are the same.
The identity bases are not.

This is a core contribution of the paper.
Reference kinds supply carriers.
Transformations determine regimes.
Counting reference kinds gives six.
The transformation analysis forces at least nine identity regimes.

% !TeX root = ../se200_referential_regimes.tex

% =============================================================================
\section{Related Work}
\label{sec:related}
% =============================================================================

This paper is situated at the intersection of formal ontology, provenance, and
classification under disagreement.
It does not propose a new upper ontology.
It derives a lower bound on identity regimes required when a substrate
is optimized for neutral reference under persistent disagreement.

% -----------------------------------------------------------------------------
\subsection{Classification under Disagreement}
\label{subsec:related_classification_disagreement}
% -----------------------------------------------------------------------------

Classification systems embed institutional, social, and evaluative commitments
rather than functioning as neutral descriptions~\citep{bowkerstar1999,haslanger2012}.
In scientific and policy settings, disagreement is often persistent rather than transient,
reflecting incompatible evidentiary, methodological, or evaluative frameworks
rather than temporary empirical uncertainty~\citep{longino1990}.
This motivates representations that remain usable without presupposing interpretive convergence.

This paper adopts that motivation but addresses a narrower structural question:
which identity distinctions become unavoidable when disagreement
is persistent and causal or normative commitment is externalized from the substrate.
The answer is conditional on the neutral-substrate constraint developed
in \emph{Neutral Substrates}~\citep{case2026neutral}
and on the formal setting fixed above.
Under those assumptions, stable reference requires at least the nine core identity regimes derived here.

% -----------------------------------------------------------------------------
\subsection{Identity, Roles, and Ontological Methodology}
\label{subsec:related_identity_roles}
% -----------------------------------------------------------------------------

Formal ontology has long emphasized that categories are individuated
by identity conditions and persistence criteria.
OntoClean clarifies why anti-rigid properties such as roles
cannot supply identity conditions~\citep{guarino1998roles,guarinowelty2002}.
Work on social and institutional entities further develops distinctions among
bearers, roles, social objects, and institutional structures~\citep{masolo2003wonderweb,guizzardi2005}.

This paper builds on those insights but applies them at the substrate level.
Its regimes are not derived from domain taxonomy.
They are derived from transformation behavior over identity bases.
The central move is to make identity operational:
a regime is determined by whether ordinary transformations
preserve, break, or leave identity neutral.

% -----------------------------------------------------------------------------
\subsection{Upper Ontologies}
\label{subsec:related_upper_ontologies}
% -----------------------------------------------------------------------------

Upper ontologies such as BFO, DOLCE, and UFO provide rich accounts of
dependence, endurants and perdurants, social objects, qualities, and
institutional entities~\citep{smithceusters2015bfo,arp2015building,gangemi2002dolce,guizzardi2005}.
They support broad ontology-engineering objectives, including reuse,
domain modeling, semantic integration, and metaphysical clarity.
The present result addresses a narrower design objective:
stable substrate-level reference under persistent disagreement.
It offers a constraint that can be imposed on
implementations when that objective is adopted.

From the perspective of BFO, several of the required carriers have natural
representational starting points.
Occurrences map naturally to occurrents.
Material plain referents map naturally to independent continuants.
Rules, records, roles, dispositions, information artifacts, and institutional
bearers can also be modeled using BFO resources, depending on the intended application.
When a BFO-based implementation is intended to serve as a
neutral substrate under persistent disagreement,
the relevant identity bases can be specified as part of the implementation.
For example, non-spatial normative applicability, record identity, and obligation-bearing capacity
may require dedicated modeling patterns to ensure that the
regime-level distinctions are preserved across transformations.

DOLCE and UFO provide rich resources for social objects,
institutional entities, qualities, roles, and dependence relations.
These resources make them natural candidates for representing
obligation-bearers, rules, and institutional structures.
Here again, the contribution is an additional constraint
imposed when neutrality under persistent
disagreement is adopted as a design objective.
A deployed substrate must make clear when it is
tracking a rule by content rather than structure,
a scope by extension rather than internal organization,
a plain referent by locus rather than object,
and a record as distinct from the thing recorded.
Where those distinctions are handled by
local modeling conventions, role predicates, or application logic,
the implementation may consider whether those conventions remain stable
when shared interpretation cannot be assumed.

These observations are illustrative.
The core is derived from the
transformation analysis of Sections~\ref{sec:algebra}--\ref{sec:bound},
not from any gaps in upper ontologies.
Upper ontologies can realize the regime structure derived here when
their implementations preserve the relevant identity behavior.
The benefit is a reusable constraint for cases where an ontology must support
neutral reference across divergent interpretations.

% -----------------------------------------------------------------------------
\subsection{Provenance, Versioning, and Exchange}
\label{subsec:related_provenance_exchange}
% -----------------------------------------------------------------------------

The transformation basis used in this paper mirrors operations
already tracked by provenance and versioning systems.
PROV represents derivation, attribution, generation, use, and revision
in ways that overlap naturally with re-expression, annotation,
refinement, branching, and provenance extension~\citep{provdm2013}.
This paper adds an identity-level constraint for cases
where provenance and versioning records
must remain referable across incompatible interpretations.

The neutral substrate and its regimes can be stated as a conservative theory
over an exchange formalism such as Common Logic~\citep{isoiec24707}.
Conservativity matters because extensions should add interpretive structure
without changing consequences about substrate-level terms~\citep{hodges1997}.
This supports heterogeneous ontology exchange and pluralistic structuring
in settings where no single interpretive framework can govern all use~\citep{kutz2010}.

Large-scale information infrastructures and digital-twin systems face
related problems of long-lived reference, versioning, and cross-institutional
reuse~\citep{edwards2011,grieves2016}.
The literature emphasizes that infrastructure must remain usable
across evolving communities and interpretations.
This paper provides a constraint that enables stable reference under persistent disagreement.

% -----------------------------------------------------------------------------
\subsection{Causal and Normative Separation}
\label{subsec:related_causal_normative}
% -----------------------------------------------------------------------------

The distinction between descriptive record and causal claim follows a familiar discipline in causal analysis.
Causal claims require additional modeling commitments and should not be confused
with descriptive structure alone~\citep{pearl2009}.
The neutral-substrate constraint applies the same discipline at the substrate boundary.
A substrate may record that a source asserted a causal or normative proposition.
It must not thereby assert the proposition as substrate-level fact.

The regimes derived here make that separation usable under transformation.
Records can persist, be annotated, forked, and extended with provenance
without becoming causal explanations or normative judgments.
Rules can remain referable without the substrate deciding which interpretation governs.
Occurrences can remain referable without the substrate deciding
their causal effects or compliance status.
That separation is essential when identity must remain stable
before interpretation is settled.

% -----------------------------------------------------------------------------
\subsection{FAIR Reuse under Persistent Disagreement}
\label{subsec:related_fair}
% -----------------------------------------------------------------------------

The FAIR principles promote data reuse across communities and over time~\citep{wilkinson2016}.
They do not specify the ontological identity conditions required
when reuse occurs under persistent legal, institutional, or analytic disagreement.
In practice, reuse is often pursued through negotiated
vocabularies, shared schemas, or domain-specific interpretive conventions.
Those strategies are valuable, but they presuppose forms of
convergence that may be unavailable in contested settings.

This result complements FAIR-aligned reuse in contested settings.
Stable reference must be maintained independently of causal or normative agreement.
The nine core regimes name identity distinctions that a neutral substrate
must represent when reuse is expected to continue across divergent interpretations.

% !TeX root = ../se200_identity_regimes.tex

% =============================================================================
\section{Limits and Boundary Conditions}
\label{sec:limits}
% =============================================================================

The result of this paper is deliberately narrow.
It applies only to neutral accountability substrates satisfying the
assumptions of Section~\ref{sec:formal_setting}.
It establishes a necessary lower bound on identity regimes:
a floor, not a ceiling.
It does not certify that a deployed system has chosen the correct regime for every real-world referent.

% -----------------------------------------------------------------------------
\subsection{Conditional Scope}
\label{subsec:limits_conditional_scope}
% -----------------------------------------------------------------------------

The lower bound depends on the neutral-substrate constraint developed in
\emph{Neutral Substrates}~\citep{case2026neutral} and on the formal setting fixed in this paper.
In particular, it assumes persistent disagreement, strong stable reference, minimality,
and exclusion of hidden regimes.
If those assumptions are rejected, the bound need not apply.

The result is structural rather than domain-specific.
It concerns identity behavior under transformation.
It does not depend on any particular legal theory,
policy framework, institutional vocabulary, or causal model.
It says that, under the stated assumptions, fewer than nine core regimes
cannot preserve stable reference for the carrier-basis pairs analyzed here
without either changing identity behavior or hiding a regime distinction elsewhere.

% -----------------------------------------------------------------------------
\subsection{No Upper Bound or Completeness Claim}
\label{subsec:limits_no_completeness}
% -----------------------------------------------------------------------------

The nine-regime lower-bound core is a lower bound only.
It leaves three questions open:
whether these nine regimes exhaust the coherent identity regimes;
whether the transformation basis $\Fset$ is complete for all possible record systems;
and whether some domains require additional reference kinds.

The monotonicity result of Corollary~\ref{se200.cor.Monotonicity} is one-directional.
Adding transformations cannot collapse distinctions already witnessed by $\Fset$.
A richer basis may leave the bound unchanged or force additional splits.
It cannot make the nine-regime lower bound disappear.

The inventory of Section~\ref{sec:inventory} is likewise scoped.
It identifies reference kinds required for the accountability settings considered here.
Additional domains may require additional carriers.
Additional carriers may require additional regimes.
They do not remove the regimes already forced for the carrier-basis pairs analyzed in this paper.

The paper does not claim that obligation-bearers, occurrences, or records
admit no further admissible identity bases.
For example, a record may be tracked by descriptive content or by a
particular carrier, file, signature, or chain-of-custody token.
An occurrence may be tracked as a fine-grained realization or as a coarser event complex.
An obligation-bearer may be tracked as a legal person, office-holder,
delegated role, or responsible organizational unit.
If such bases are required by a deployment, they may force additional regimes.
That possibility does not weaken the lower bound.
It confirms that the present result is a lower-bound core, not a completeness theorem.

% -----------------------------------------------------------------------------
\subsection{Excluded Systems}
\label{subsec:limits_excluded_systems}
% -----------------------------------------------------------------------------

Several classes of systems fall outside the scope of the result.

\paragraph{Interpretively dependent identity.}
A system is outside the scope if the identity condition of a referent depends
on causal, normative, evaluative, or methodological commitments supplied by an extension.
Such a system may be useful within a coordinated framework.
It is not a neutral substrate in the sense analyzed here.

\paragraph{Relation-determined identity.}
A system is outside the scope if relations introduce, modify, or
determine the identity conditions of their relata.
Relations may connect already-individuated referents.
They may not supply the identity basis that makes those referents the same or different.

\paragraph{Extension-variant identity.}
A system is outside the scope if identity conditions vary when admissible extensions are added.
Such systems violate the stable-reference requirement.
They may support local or negotiated identity.
They do not support the strong form of stable reference studied here.

\paragraph{Embedded interpretive commitments.}
A system is outside the scope if it makes substrate-level causal or normative commitments
whose validity depends on a particular interpretive framework.
The substrate may record attributed causal or normative claims.
It may not assert them as object-level substrate facts.

\paragraph{Undeclared hidden regimes.}
A system is outside the scope if it collapses regimes
and then recovers the lost distinction through roles, flags,
contextual predicates, workflow states, schema conventions, or application-level discriminators.
Such a system has not reduced structure.
It has relocated structure into an interpretive mechanism.

\paragraph{Unmodeled transformation behavior.}
A system is outside the scope when its essential identity behavior
depends on transformations not represented in $\Fset$ and not reducible to the
transformation families used here.
Such a case does not refute the lower bound.
It indicates that the transformation basis should be enlarged.
By monotonicity, such enlargement may preserve the bound or force additional regimes.

% -----------------------------------------------------------------------------
\subsection{Non-Goals and Weakened Objectives}
\label{subsec:limits_nongoals}
% -----------------------------------------------------------------------------

Several questions remain outside this paper.
This paper constructs no causal or normative extensions.
It takes no position on which causal account is correct or which normative interpretation should govern.
It offers no domain-specific modeling guidance.
It does not claim that all useful ontologies should be neutral substrates.
It does not claim that hidden regimes are always poor engineering.
Its claim is only that hidden regimes are incompatible with the specific objective analyzed here:
stable substrate-level reference under persistent disagreement.

If stability is weakened to permit context-dependent regime membership,
some distinctions may be represented as role-relative or parameter-dependent refinements of fewer base types.
That may be a reasonable engineering choice in coordinated settings.
But then the structural burden moves into coordination, versioning, governance,
or application logic whose admissibility is not fixed at the substrate level.
The lower bound proved here does not apply under that weakened objective.

% -----------------------------------------------------------------------------
\subsection{Boundary with Accountable Records}
\label{subsec:limits_boundary_accountable_records}
% -----------------------------------------------------------------------------

This paper concerns identity.
It derives the lower-bound regime structure required for stable reference.
It does not provide a conformance checker for deployed record systems.
It does not specify how schemas must declare bases, logs, or discriminator surfaces.
It does not decide whether an implementation has disclosed every identity-relevant surface.

Those questions are deferred to the companion \emph{Accountable Records} paper.
That paper is intended to develop accountable records, basis declaration,
transformation logs, and conformance checks.
The present paper supplies the lower-bound identity core that such a conformance analysis requires.
It establishes part of the identity cost of the neutral-substrate constraint,
while \emph{Accountable Records} studies how that cost
can be made checkable in deployed record systems.

% -----------------------------------------------------------------------------
\subsection{Boundary of the Claim}
\label{subsec:limits_boundary_claim}
% -----------------------------------------------------------------------------

Within the stated scope, the conclusion is unavoidable.
A neutral accountability substrate that supports the inventory of Section~\ref{sec:inventory},
preserves stable reference under the transformation basis $\Fset$,
and excludes hidden regimes must realize
at least the nine identity regimes in the lower-bound core.
Outside that scope, different optimization objectives may yield different structures.
The result is best understood as a constraint theorem:
for neutral substrate-level reference under persistent disagreement,
at least this much identity structure is required.

% !TeX root = ../se200_referential_regimes.tex

% =============================================================================
\section{Conclusion}
\label{sec:conclusion}
% =============================================================================

Identity in a neutral substrate is not fixed by naming kinds alone.
It is fixed by specifying which transformations preserve sameness.
Once the ordinary operations of accountability systems are specified,
individuation becomes a structural question rather than a taxonomic preference.

This paper began with the neutral-substrate constraint developed in
\emph{Neutral Substrates}~\citep{case2026neutral}.
A substrate that must remain usable under persistent disagreement
may commit to stable reference and permitted attribution,
but not to object-level causal or normative interpretation.
This paper explores what identity structure that constraint requires
when records are transformed in ordinary use.

The answer is a lower-bound core.
An accountability substrate must be able to refer to six kinds of things:
obligation-bearers, rules, occurrences, scopes, records, and plain referents.
Those six reference kinds supply carriers, but do not settle identity.
Three of them are assigned more than one admissible basis:
a scope by extension or by structure,
a rule by content or by structure,
a plain referent by locus or by object.
Ordinary transformations force apart those alternatives:
decomposition separates extension-fixed from structure-fixed scopes,
refinement separates content-fixed from structure-fixed rules,
and forking separates locus-fixed from object-fixed plain referents.
Thus six reference kinds force at least nine identity regimes:
\[
\begin{aligned}
&\OBL,\quad \OCC,\quad \REC,\quad \LOC,\quad \OBJ,\\
&\SCOPEE,\quad \SCOPES,\quad \RULEC,\quad \RULES .
\end{aligned}
\]

The central result is the identification of the additional regimes.
They are not additional reference kinds, but distinct identity bases forced apart by transformation behavior.
An analysis that names kinds may undercount the structure required for stable reference.
Collapsing regimes does not produce genuine neutrality.
A collapsed distinction must change identity behavior
or be recovered through a hidden discriminator.

The result is conditional.
It holds under the assumptions fixed in this paper:
persistent disagreement, strong stable reference, minimality,
and exclusion of hidden regimes.
It is a lower bound rather than an upper bound or a completeness claim.
A richer transformation basis may leave it unchanged or force additional regimes.

The result is diagnostic rather than prescriptive.
It does not claim that every ontology should expose these nine regimes,
nor that contextual typing, role predicates, or hidden discriminators
are useless as engineering devices.
It says that such devices cannot substitute for regime-level identity commitments
in a neutral substrate under persistent disagreement.
When stable reference must be held as common ground before
interpretation is settled, the substrate structure must preserve
the relevant identity distinctions.

Three papers in this series explore the implications of that constraint.
\emph{Neutral Substrates} states the neutrality constraint.
\emph{Referential Regimes} measures its identity cost.
\emph{Accountable Records} studies how that cost can be made checkable in deployed record systems.
This work addresses the middle step:
stable reference under persistent disagreement requires
transformation-invariant identity regimes,
and at least nine such regimes are forced by the ordinary operations
accountability systems already perform.


% =============================================================================
\section*{Statements and Declarations}
% =============================================================================

\subsection*{Acknowledgments} The author thanks the anonymous reviewers
of an earlier version of this work for comments that helped improve
the framing, organization, and presentation of the work.

\subsection*{Author Contributions}
The author is the sole contributor to this work and
is responsible for all aspects of the research, authorship, and publication.

\subsection*{Use of AI-Assisted Tools}
AI-assisted tools were used for editing, formatting, and consistency checking.
The author reviewed all suggestions and is solely responsible for the content.

\subsection*{Declaration of Conflicting Interest}
The author declares no potential conflicts of interest with respect to the
research, authorship, and publication of this work.

\appendix
\clearpage
% !TeX root = ../se200_identity_regimes.tex

% =============================================================================
\section{Pairwise Separation}
\label{app:sep}
% =============================================================================

The nine regimes yield $\binom{9}{2}=36$ unordered pairs.
Each pair is separated either by carrier difference or, for same-carrier split pairs,
by the canonical operation that forces the split.
Only three pairs share a reference kind and differ by identity basis:
$\LOC/\OBJ$, $\SCOPEE/\SCOPES$, and $\RULEC/\RULES$.
All other pairs are separated by carrier difference.

\begin{center}
\small
\setlength{\tabcolsep}{2pt}
\begin{tabular}{@{}l ccccccccc@{}}
\toprule
 & $\OBL$ & $\OCC$ & $\REC$ & $\LOC$ & $\OBJ$ & $\SCOPEE$ & $\SCOPES$ & $\RULEC$ & $\RULES$\\
\midrule
$\OBL$   & -- & $\bullet$ & $\bullet$ & $\bullet$ & $\bullet$ & $\bullet$ & $\bullet$ & $\bullet$ & $\bullet$\\
$\OCC$   & $\bullet$ & -- & $\bullet$ & $\bullet$ & $\bullet$ & $\bullet$ & $\bullet$ & $\bullet$ & $\bullet$\\
$\REC$   & $\bullet$ & $\bullet$ & -- & $\bullet$ & $\bullet$ & $\bullet$ & $\bullet$ & $\bullet$ & $\bullet$\\
$\LOC$  & $\bullet$ & $\bullet$ & $\bullet$ & -- & BF & $\bullet$ & $\bullet$ & $\bullet$ & $\bullet$\\
$\OBJ$  & $\bullet$ & $\bullet$ & $\bullet$ & BF & -- & $\bullet$ & $\bullet$ & $\bullet$ & $\bullet$\\
$\SCOPEE$  & $\bullet$ & $\bullet$ & $\bullet$ & $\bullet$ & $\bullet$ & -- & AD & $\bullet$ & $\bullet$\\
$\SCOPES$  & $\bullet$ & $\bullet$ & $\bullet$ & $\bullet$ & $\bullet$ & AD & -- & $\bullet$ & $\bullet$\\
$\RULEC$  & $\bullet$ & $\bullet$ & $\bullet$ & $\bullet$ & $\bullet$ & $\bullet$ & $\bullet$ & -- & RF\\
$\RULES$  & $\bullet$ & $\bullet$ & $\bullet$ & $\bullet$ & $\bullet$ & $\bullet$ & $\bullet$ & RF & --\\
\bottomrule
\end{tabular}
\end{center}

Here $\bullet$ indicates separation by carrier difference.
The entries BF, AD, and RF indicate same-kind split pairs separated by the corresponding canonical operation.

The $\LOC/\OBJ$ pair is separated by branch/fork.
A fork over a fixed locus may preserve locus identity while breaking object identity.
Thus the two regimes induce distinct identity relations.

The $\SCOPEE/\SCOPES$ pair is separated by aggregation/decomposition.
A decomposition may preserve the covered cases while changing the internal scope structure.
Thus the two regimes induce distinct identity relations.

The $\RULEC/\RULES$ pair is separated by refinement.
A refinement may preserve rule content while changing textual or internal rule structure.
Thus the two regimes induce distinct identity relations.

Therefore every pair of regimes is separated either by carrier difference or by a transformation witness.
No pair collapses.
   % pairwise collapse enumeration

\begin{thebibliography}{99}
  \small

  \bibitem[Arp et al.(2015)]{arp2015building}
  R.\ Arp, B.\ Smith, and A.\ D.\ Spear.
  \emph{Building Ontologies with Basic Formal Ontology}.
  MIT Press, 2015.

  \bibitem[Bowker and Star(1999)]{bowkerstar1999}
  G.\ C.\ Bowker and S.\ L.\ Star.
  \emph{Sorting Things Out: Classification and Its Consequences}.
  MIT Press, Cambridge, MA, 1999.

  \bibitem[Case(2026)]{case2026neutral}
  D.\ M.\ Case.
  \emph{Neutral Substrates: A Design Constraint for Interoperability under Persistent Disagreement}.
  arXiv preprint \href{https://arxiv.org/abs/2601.14271}{arXiv:2601.14271}, 2026.

  \bibitem[Edwards(2011)]{edwards2011}
  P.\ N.\ Edwards.
  \emph{A Vast Machine: Computer Models, Climate Data, and the Politics of Global Warming}.
  MIT Press, 2011.

  \bibitem[Gangemi et al.(2002)]{gangemi2002dolce}
  A.\ Gangemi, N.\ Guarino, C.\ Masolo, A.\ Oltramari, and L.\ Schneider.
  Sweetening ontologies with DOLCE.
  In \emph{EKAW 2002}, volume 2473 of LNCS, pages 166--181. Springer, 2002.

  \bibitem[Grieves and Vickers(2016)]{grieves2016}
  M.\ Grieves and J.\ Vickers.
  Digital twin: Mitigating unpredictable, undesirable emergent behavior in complex systems.
  In \emph{Transdisciplinary Perspectives on Complex Systems}, pages 85--113. Springer, 2016.

  \bibitem[Guarino(1998)]{guarino1998roles}
  N.\ Guarino.
  On the ontological status of roles.
  \emph{Applied Ontology}, 1998.

  \bibitem[Guarino and Welty(2002)]{guarinowelty2002}
  N.\ Guarino and C.\ A.\ Welty.
  Evaluating ontological decisions with OntoClean.
  \emph{Communications of the ACM}, 45(2):61--65, 2002.

  \bibitem[Guizzardi(2005)]{guizzardi2005}
  G.\ Guizzardi.
  \emph{Ontological Foundations for Structural Conceptual Models}.
  PhD thesis, University of Twente, 2005.

  \bibitem[Haslanger(2012)]{haslanger2012}
  S.\ Haslanger.
  \emph{Resisting Reality: Social Construction and Social Critique}.
  Oxford University Press, 2012.

  \bibitem[Hodges(1997)]{hodges1997}
  W.\ Hodges.
  \emph{A Shorter Model Theory}.
  Cambridge University Press, 1997.

  \bibitem[ISO/IEC(2018)]{isoiec24707}
  ISO/IEC.
  Information technology --- Common Logic (CL): A framework for a family of logic-based languages.
  ISO/IEC 24707:2018, 2018.

  \bibitem[Kutz et al.(2010)]{kutz2010}
  O.\ Kutz, T.\ Mossakowski, and D.\ L\"ucke.
  Carnap, Goguen, and the hyperontologies: Logical pluralism and heterogeneous structuring in ontology design.
  \emph{Logica Universalis}, 4(2):255--333, 2010.

  \bibitem[Longino(1990)]{longino1990}
  H.\ E.\ Longino.
  \emph{Science as Social Knowledge: Values and Objectivity in Scientific Inquiry}.
  Princeton University Press, 1990.

  \bibitem[Masolo et al.(2003)]{masolo2003wonderweb}
  C.\ Masolo, S.\ Borgo, A.\ Gangemi, N.\ Guarino, and A.\ Oltramari.
  WonderWeb Deliverable D18: Ontology Library.
  IST Project 2001-33052 WonderWeb, 2003.

  \bibitem[Moreau and Missier(2013)]{provdm2013}
  L.\ Moreau and P.\ Missier (eds.).
  PROV-DM: The PROV Data Model.
  W3C Recommendation, 30 April 2013.

  \bibitem[Pearl(2009)]{pearl2009}
  J.\ Pearl.
  \emph{Causality: Models, Reasoning, and Inference}.
  Cambridge University Press, 2nd edition, 2009.

  \bibitem[Smith and Ceusters(2015)]{smithceusters2015bfo}
  B.\ Smith and W.\ Ceusters.
  Basic Formal Ontology 2.0: Specification and user's guide.
  \emph{Applied Ontology}, 10(1):1--77, 2015.

  \bibitem[Wilkinson et al.(2016)]{wilkinson2016}
  M.\ D.\ Wilkinson et al.
  The FAIR Guiding Principles for scientific data management and stewardship.
  \emph{Scientific Data}, 3:160018, 2016.

\end{thebibliography}

\end{document}
